Lattice summation and Fourier transform supply arithmetic relations before
any spectral assumption is made. The first question is whether these
relations respect the complete Weil functional. A positive operator
realization asks a further question: can the quotient carry a positive
Hilbert representation whose ordinary trace is that functional?
We construct the quotient and its scalar functional, then determine the
form of every trace-class representation that could realize it.

\section{The test algebra and its scalar functional}

All coefficient spaces are complex. The additive Haar measure is $du$,
and Hilbert inner products are conjugate-linear in the first variable.
For $a>1/2$ put
\begin{equation}\label{nwc:eq:tests}
\begin{gathered}
 \mathcal A_a=\{g\in C^\infty(\mathbb R):p_{a,N}(g)<\infty
       \text{ for every }N\geq0\},\\
 p_{a,N}(g)=\sum_{j=0}^N\int_{\mathbb R}e^{a|u|}|g^{(j)}(u)|\,du.
\end{gathered}
\end{equation}
The same formula with $a=0$ defines the unweighted seminorm $p_{0,N}$.
Use convolution and the involution
\[
 (f*g)(u)=\int_{\mathbb R}f(t)g(u-t)\,dt,
 \qquad g^*(u)=\overline{g(-u)}.
\]
Let $\mathcal A=\bigcup_{n\geq1}\mathcal A_{1/2+1/n}$ with the final
locally convex topology. Define
\begin{equation}\label{nwc:eq:transforms}
 h_g(z)=\int g(u)e^{izu}\,du,
 \qquad H_g(s)=\int g(u)e^{(s-1/2)u}\,du
             =h_g(-i(s-1/2)).
\end{equation}
Thus $H_g(1/2+ir)=h_g(r)$. No integrable test is the identity for
convolution; these algebras are nonunital.

\begin{lemma}
Each $\mathcal A_a$\label{nwc:lem:tests} is a complete Fr\'echet star algebra, and
$C_c^\infty(\mathbb R)$ is dense in it. Its seminorms satisfy
\[
 p_{a,N}(f*g)\leq p_{a,N}(f)p_{a,0}(g),\qquad
 p_{a,N}(g^*)=p_{a,N}(g).
\]
For $0\leq b<a$ and integers $m,k\geq0$,
\begin{equation}\label{nwc:eq:strip}
 \sup_{|\operatorname{Im}z|\leq b}
 (1+|\operatorname{Re}z|)^m|h_g^{(k)}(z)|
 \leq C_{a,b,m,k}p_{a,m}(g).
\end{equation}
The transform is holomorphic on $|\operatorname{Im}z|<a$ and satisfies
$h_{f*g}=h_fh_g$ and $h_{g^*}(z)=\overline{h_g(\bar z)}$.
\end{lemma}
\begin{proof}
The weight is submultiplicative, so Fubini gives the convolution
inequality. A Cauchy sequence has weighted $L^1$ limits of all its
derivatives. Distributional differentiation identifies successive
limits; local Sobolev regularity then gives a common smooth
representative. This proves completeness.
If $\chi$ is a smooth compact cutoff equal to one near zero,
Leibniz's rule and integrable tails give $\chi(u/R)g(u)\to g(u)$ in
every seminorm.

Differentiating the transform $k$ times inserts $(iu)^k$.
Integration by parts $m$ times gives, on the smaller strip,
\[
 |z|^m|h_g^{(k)}(z)|\leq
 \sum_{j=0}^{\min(m,k)}\binom mj\frac{k!}{(k-j)!}
 \int |u|^{k-j}|g^{(m-j)}(u)|e^{b|u|}\,du.
\]
Each polynomial is absorbed by $e^{(a-b)|u|}$.
The boundary terms vanish because the relevant weighted functions and
their first derivatives are integrable. The direct integral estimate
handles $|z|\leq1$. Dominated differentiation and Fubini prove the
remaining assertions.
\end{proof}

Write $\Lambda(p^k)=\log p$ for primes $p$ and integers $k\geq1$,
and $\Lambda(n)=0$ otherwise. Let $\psi=\Gamma'/\Gamma$.
The scalar functional used throughout is
\begin{equation}\label{nwc:eq:weil}
\begin{aligned}
 L(g)={}&H_g(0)+H_g(1)-g(0)\log\pi\\
 &+\frac1{2\pi}\int_{\mathbb R}h_g(r)
                  \operatorname{Re}\psi(1/4+ir/2)\,dr\\
 &-\sum_{n\geq2}\frac{\Lambda(n)}{\sqrt n}
                 \bigl(g(\log n)+g(-\log n)\bigr).
\end{aligned}
\end{equation}
Both polar evaluations and both orientations of every prime power
belong to this definition. These conventions agree with
\cite{nwc:CC}*{Appendices A--B} after the substitution
$k(x)=x^{-1/2}g(\log x)$.

\begin{lemma}\label{nwc:lem:weil-continuity}
All terms in \eqref{nwc:eq:weil} converge absolutely, and
$|L(g)|\leq C_a p_{a,2}(g)$ on every test step.
Moreover $L(g^*)=\overline{L(g)}$.
\end{lemma}
\begin{proof}
The function $e^{a|u|}g(u)$ and its weak derivative are integrable,
so it vanishes at both ends and
\[
 \sup_u e^{a|u|}|g(u)|\leq
 a\|e^{a|u|}g\|_1+\|e^{a|u|}g'\|_1.
\]
The prime terms are consequently bounded by
$2(a+1)p_{a,1}(g)\sum_{n\geq2}(\log n)n^{-a-1/2}$.
The poles are bounded by $2p_{a,0}(g)$.
The logarithmic bound
$\psi(1/4+ir/2)=O(1+\log(2+|r|))$, proved below, and
\eqref{nwc:eq:strip} with $m=2$ control the integral.
Reflection and conjugation in the displayed formula give the last
identity.
\end{proof}

\section{Lattice relations and their topology}

Let $\mathcal S(\mathbb R)$ have seminorms
$q_{M,k}(v)=\sup_x(1+|x|)^M|v^{(k)}(x)|$, for nonnegative integers
$M,k$. The closed source subspace is
\[
 \mathcal S_{\mathrm{ev},0}
 =\{v\in\mathcal S(\mathbb R):v(-x)=v(x),\ v(0)=0,\ \int v=0\}.
\]
Define Fourier transform, normalized dilation, and lattice summation by
\begin{equation}\label{nwc:eq:theta}
 \mathcal Fv(y)=\int v(x)e^{-2\pi ixy}\,dx,\qquad
 V_tv(x)=e^{-t/2}v(e^{-t}x),\qquad
 \mathcal Ev(u)=e^{u/2}\sum_{n\geq1}v(ne^u).
\end{equation}
The lattice sum has coefficient one and runs over positive integers.
Write $\zeta(s)=\sum_{n\geq1}n^{-s}$ for $\operatorname{Re}s>1$.

\begin{theorem}\label{nwc:thm:theta}
The map $\mathcal E$ is continuous from $\mathcal S_{\mathrm{ev},0}$
to every $\mathcal A_a$. If $R_tg(u)=g(u-t)$ and
$W=1/2+x\partial_x$, then
\begin{equation}\label{nwc:eq:intertwining}
 \mathcal EV_t=R_t\mathcal E,\quad
 \mathcal E\mathcal Fv(u)=\mathcal Ev(-u),\quad
 (\mathcal Ev)^*=\mathcal E(\mathcal F\bar v),\quad
 \partial_u^j\mathcal Ev=\mathcal E(W^jv).
\end{equation}
To give an explicit bound, write
\[
 W^j=\sum_{k=0}^j c_{jk}x^k\partial_x^k,
 \quad c_{00}=1,\quad
 c_{j+1,k}=(k+1/2)c_{jk}+c_{j,k-1},
 \quad Q_{M,j}(v)=\sum_{k=0}^jc_{jk}q_{M+k,k}(v),
\]
with coefficients outside their ranges equal to zero.
For every integer $M>\max(1,a+1/2)$,
\begin{equation}\label{nwc:eq:theta-bound}
 p_{a,N}(\mathcal Ev)\leq
 \frac{\zeta(M)}{M-a-1/2}
 \sum_{j=0}^N\bigl(Q_{M,j}(v)+Q_{M,j}(\mathcal Fv)\bigr).
\end{equation}
\end{theorem}
\begin{proof}
Fourier differentiation and integration by parts give
\[
 q_{M,k}(\mathcal Fv)\leq
 2\sum_{r=0}^M\binom Mr(2\pi)^{k-r}
 \sum_{j=0}^{\min(r,k)}\binom rj\frac{k!}{(k-j)!}
 q_{k-j+2,r-j}(v).
\]
The factor two is $\int_{\mathbb R}(1+|x|)^{-2}\,dx$.
Thus Fourier transform is continuous on Schwartz space.
Fourier inversion follows by inserting $e^{-\pi\epsilon y^2}$
in the inverse integral and using the Gaussian approximate identity.
It gives $\mathcal F^2v(x)=v(-x)$, and shows that $\mathcal F$
preserves the two vanishing conditions.

For $y>0$, the smooth periodization
$\sum_{n\in\mathbb Z}v(y(t+n))$ has Fourier coefficients
$y^{-1}\mathcal Fv(k/y)$. Both series converge normally with all
derivatives. Evaluation at $t=0$ gives Poisson summation.
The zero terms vanish, and evenness gives
\begin{equation}\label{nwc:eq:poisson}
 \sum_{n\geq1}v(ny)=y^{-1}\sum_{n\geq1}\mathcal Fv(n/y).
\end{equation}
This proves the reflection identity in \eqref{nwc:eq:intertwining}.
The normalization is also that of \cite{nwc:Serre}*{VII, \S6}.

The theta series can be differentiated on every compact $u$ interval.
Differentiation gives $\partial_u^j\mathcal Ev=\mathcal E(W^jv)$;
expanding $W^j$ gives the stated nonnegative coefficients.
For $u\geq0$,
\[
 |\partial_u^j\mathcal Ev(u)|
 \leq\zeta(M)e^{-(M-1/2)u}Q_{M,j}(v).
\]
For $u\leq0$, \eqref{nwc:eq:poisson} gives the same bound with
$u$ replaced by $-u$ and $v$ by $\mathcal Fv$.
Integration proves \eqref{nwc:eq:theta-bound}.

A change of variables gives
$\mathcal FV_t=V_{-t}\mathcal F$ and $\mathcal EV_t=R_t\mathcal E$.
Conjugation commutes with $\mathcal E$, which proves its star identity.
The operator $W$ preserves the source: $(Wv)(0)=v(0)/2$ and
$\int Wv=-(1/2)\int v$. Hence every identity has the stated domain.
\end{proof}

For a general even Schwartz function the omitted lattice zero contributes
\[
 \mathcal Ev(u)=\mathcal E\mathcal Fv(-u)
     +\tfrac12e^{-u/2}\mathcal Fv(0)-\tfrac12e^{u/2}v(0).
\]
The two vanishing conditions remove these separate tails.
Dilation and translation have the bounds
\begin{equation}\label{nwc:eq:translation-bounds}
 \begin{split}
 q_{M,k}(V_tv)&\leq e^{-(k+1/2)t}\max(1,e^t)^M q_{M,k}(v),\\
 p_{a,N}(R_tg)&\leq e^{a|t|}p_{a,N}(g),\qquad
 p_{a,N}(g')\leq p_{a,N+1}(g).
 \end{split}
\end{equation}
They form smooth groups, with generators $-W$ and $-\partial_u$.
In particular, dilation by two composes with itself to dilation by four
already on the source:
$V_{\log2}^2=V_{\log4}$ and $R_{\log2}^2=R_{\log4}$.

\begin{proposition}\label{nwc:prop:ideal}
The closed convolution star ideal generated by the theta images is
\[
 \mathcal J_a=\overline{\mathcal E(\mathcal S_{\mathrm{ev},0})}^{\,\mathcal A_a}.
\]
It is stable under reflection, conjugation, $R_t$, and $\partial_u$.
The quotient $\mathcal Q_{\Theta,a}=\mathcal A_a/\mathcal J_a$ is a
complete Hausdorff Fr\'echet star algebra.
For $b>a$ there is a continuous dense-range map
\begin{equation}\label{nwc:eq:step-map}
 \mathcal Q_{\Theta,b}\longrightarrow\mathcal Q_{\Theta,a},
 \qquad
 \ker= (\mathcal A_b\cap\mathcal J_a)/\mathcal J_b.
\end{equation}
\end{proposition}
\begin{proof}
The theta image is a linear star-stable subspace.
For $f\in C_c^\infty(\mathbb R)$ the integral
$v_f=\int f(t)V_tv\,dt$ converges in Schwartz space, by
\eqref{nwc:eq:translation-bounds}, and retains the two vanishing
conditions. Continuity of $\mathcal E$ gives
$\mathcal Ev_f=f*\mathcal Ev$.
Approximation of an arbitrary $f\in\mathcal A_a$ by compact tests,
then closure in the second factor, proves ideal stability.
The source identities and their bounds prove the other invariances.
The quotient is complete because the ideal is closed.

The inclusion $\mathcal A_b\to\mathcal A_a$ has seminorm bound one
and sends $\mathcal J_b$ into $\mathcal J_a$.
Its induced quotient map has the displayed kernel.
The compact core is contained in $\mathcal A_b$, which gives dense range.
\end{proof}

Injectivity in \eqref{nwc:eq:step-map} would require
$\mathcal A_b\cap\mathcal J_a=\mathcal J_b$.
No such equality is needed below. The quotient maps form a compatible
system; its locally convex direct limit is not asserted to be strict
or Hausdorff.

\section{The Mellin factor and scalar descent}

The continuation of $\zeta$ used below follows from the integral formulas in the
proof of Theorem~\ref{nwc:thm:descent}.
For $v\in\mathcal S_{\mathrm{ev},0}$ set
$M_v(s)=\int_0^\infty v(x)x^{s-1}\,dx$.
Evenness and $v(0)=0$ give $v(x)=O(x^2)$ at zero.
Thus $M_v$ is holomorphic on $\operatorname{Re}s>-2$.
On a closed strip $\sigma_-\leq\operatorname{Re}s\leq\sigma_+$
in that half-plane, choose a nonnegative integer $B>\sigma_+$.
For $k\geq0$,
\[
 |M_v^{(k)}(s)|\leq
 \frac{k!q_{0,2}(v)}{2(\sigma_-+2)^{k+1}}
 +\frac{k!q_{B,0}(v)}{(B-\sigma_+)^{k+1}}.
\]
Integration by parts gives
$s^mM_v(s)=(-1)^mM_{(x\partial_x)^mv}(s)$.
These bounds prove rapid vertical decrease, with every derivative,
uniformly on each such strip.

\begin{lemma}\label{nwc:lem:mellin}
The theta map is injective and satisfies
\begin{equation}\label{nwc:eq:mellin-factor}
 H_{\mathcal Ev}(s)=\zeta(s)M_v(s),\qquad
 h_{\mathcal Ev}(z)=\zeta(1/2+iz)M_v(1/2+iz).
\end{equation}
The first equality initially holds for $\operatorname{Re}s>1$ and
continues throughout $\operatorname{Re}s>-2$.
Its left side supplies an entire continuation of the product.
For every real $r$,
\begin{equation}\label{nwc:eq:real-hull}
 h_j(r)=0\text{ for all }j\in\mathcal J_a
 \quad\Longleftrightarrow\quad \zeta(1/2+ir)=0.
\end{equation}
\end{lemma}
\begin{proof}
Absolute convergence for $\operatorname{Re}s>1$ permits interchange
of sum and integral after $x=e^u$:
\[
 H_{\mathcal Ev}(s)=\int_0^\infty\sum_{n\geq1}v(nx)x^{s-1}\,dx
                  =\zeta(s)M_v(s).
\]
Since $M_v(1)=0$, the pole at one is removable.
Analytic continuation proves the asserted extension.
The theta estimate makes the left side entire.
If it is zero, the Euler product gives $M_v=0$ on
$\operatorname{Re}s>1$; Mellin inversion gives $v=0$.

Vanishing of zeta implies vanishing on the whole ideal by continuity
and multiplicativity of evaluation.
Conversely, fix $s_0=1/2+ir$ and a nonzero nonnegative
$\varphi\in C_c^\infty(1,2)$.
For $x>0$ put $w(x)=x^{-ir}\varphi(x)$ and
$v(x)=\frac12w(x/2)-w(x)$, and extend $v$ evenly.
This is an allowed source function, with
\[
 M_v(s_0)=(2^{s_0-1}-1)\int_0^\infty\varphi(x)x^{-1/2}\,dx\ne0.
\]
The first factor is nonzero because $|2^{s_0-1}|=2^{-1/2}$.
This proves \eqref{nwc:eq:real-hull}.
\end{proof}

The Mellin transform $M_v$ itself need not be entire; Taylor subtraction
at zero can give poles at the negative even integers.
Equation \eqref{nwc:eq:mellin-factor} concerns the entire continued
product. It also shows that every theta relation vanishes to at least
the zeta multiplicity at each nontrivial zero. The strip derivative
bounds extend this assertion to $\mathcal J_a$.

\begin{theorem}
For every $a>1/2$,\label{nwc:thm:descent} the full functional \eqref{nwc:eq:weil} annihilates
$\mathcal J_a$. It therefore defines compatible continuous Hermitian
functionals $\overline L_a$ on $\mathcal Q_{\Theta,a}$, with
\[
 |\overline L_a([g])|\leq C_a\inf_{j\in\mathcal J_a}p_{a,2}(g+j).
\]
This descent does not require positivity or a hypothesis on the zeros.
\end{theorem}
\begin{proof}
We first establish the analytic ingredients of a contour calculation.
Unique factorization gives, for $\operatorname{Re}s>1$,
\begin{equation}\label{nwc:eq:dirichlet}
 \frac{\zeta'}{\zeta}(s)=-\sum_{n\geq2}\Lambda(n)n^{-s}.
\end{equation}
An integration of the integer counting function followed by one
integration by parts gives
\begin{equation}\label{nwc:eq:euler-maclaurin}
 \zeta(s)=\frac1{s-1}+\frac12+\frac{s}{12}
 -\frac{s(s+1)}2\int_1^\infty B_2(\{x\})x^{-s-2}\,dx,
 \quad B_2(t)=t^2-t+\tfrac16.
\end{equation}
The periodic $B_2$ is continuous at integers, and its derivative
between integers is $2(\{x\}-1/2)$.
The integral is holomorphic on $\operatorname{Re}s>-1$.
On every closed strip there, it and its $s$ derivative are bounded.
Consequently $|\zeta(s)|+|\zeta'(s)|=O((1+|\operatorname{Im}s|)^2)$
for $|\operatorname{Im}s|\geq1$.

Let $\gamma_E=\lim_{N\to\infty}(\sum_{n=1}^N n^{-1}-\log N)$.
The beta integral gives Euler's limit for gamma, and hence its locally
normally convergent reciprocal product
\[
 \frac1{\Gamma(z)}=ze^{\gamma_Ez}
           \prod_{n\geq1}(1+z/n)e^{-z/n}.
\]
Indeed, the substitution $u=Nt$ in the beta integral gives the limit
by dominated convergence. Taking logarithmic derivatives gives
\[
 \psi(z)=-\gamma_E-\frac1z+
       \sum_{n\geq1}\left(\frac1n-\frac1{n+z}\right).
\]
On a bounded real strip with $|\operatorname{Im}z|\geq1$, split
this sum at $n$ of order $|\operatorname{Im}z|$.
The initial harmonic sum is logarithmic; the remaining tail is
$O(\sum_{n>T}|z|n^{-2})=O(1)$.
Thus $\psi(z)=O(1+\log(2+|\operatorname{Im}z|))$ there.
The same product shows that gamma has no zeros and has poles only
at the nonpositive integers.

Put $\theta(t)=\sum_{n\geq1}e^{-\pi n^2t}$.
Poisson summation for the Gaussian, then splitting the Mellin integral
at one, gives
\begin{equation}\label{nwc:eq:xi-integral}
 \pi^{-s/2}\Gamma(s/2)\zeta(s)
 =\frac1{s(s-1)}+
 \int_1^\infty\theta(t)
       \bigl(t^{s/2-1}+t^{(1-s)/2-1}\bigr)\,dt.
\end{equation}
The integral is entire by exponential decay.
It follows that
\[
 \xi(s)=\tfrac12s(s-1)\pi^{-s/2}\Gamma(s/2)\zeta(s)
\]
is entire, satisfies $\xi(s)=\xi(1-s)$, and has
$\xi(0)=\xi(1)=1/2$.
Writing $\nu=\xi'/\xi$ as a meromorphic function, we have
\begin{equation}\label{nwc:eq:logderivative}
 \nu(1-s)=-\nu(s),\qquad
 \nu(s)=\frac1s+\frac1{s-1}-\frac12\log\pi
               +\frac12\psi(s/2)+\frac{\zeta'}{\zeta}(s).
\end{equation}

Choose $1<c<\min(3/2,a+1/2)$ and let
$I_c(F)=(2\pi i)^{-1}\int_{c-i\infty}^{c+i\infty}F(s)\,ds$,
with upward orientation. For every $g\in\mathcal A_a$,
\begin{equation}\label{nwc:eq:scalar-contour}
 L(g)=I_c\bigl(\nu(s)(H_g(s)+H_g(1-s))\bigr).
\end{equation}
To verify this identity, expand \eqref{nwc:eq:logderivative}.
For $U(s)=H_g(s)+H_g(1-s)$ and $r(s)=1/s+1/(s-1)$,
shift $I_c(rU)$ to $\operatorname{Re}s=1/2$.
The residue crossed is $H_g(0)+H_g(1)$; the middle integral vanishes
because $r(1-s)=-r(s)$ and $U(1-s)=U(s)$.
Fourier inversion gives the constant contribution $-g(0)\log\pi$.
The gamma term shifts to the middle line without a pole and becomes
\[
 I_c\bigl(\tfrac12\psi(s/2)U(s)\bigr)
 =\frac1{2\pi}\int h_g(r)\operatorname{Re}\psi(1/4+ir/2)\,dr.
\]
Finally, Mellin inversion gives
\[
 I_c(n^{-s}H_g(s))=n^{-1/2}g(\log n),\qquad
 I_c(n^{-s}H_g(1-s))=n^{-1/2}g(-\log n).
\]
Absolute convergence of \eqref{nwc:eq:dirichlet} proves the prime
part with its negative sign. The strip and gamma bounds make every
horizontal integral tend to zero.

Now let $g=\mathcal Ev$. On $1-c\leq\operatorname{Re}s\leq c$,
the function $\Phi(s)=\nu(s)H_g(s)$ extends holomorphically.
Away from zero and one the gamma and polynomial factors in $\xi$
are nonzero, so every logarithmic pole at a zero cancels against
the zeta factor in \eqref{nwc:eq:mellin-factor}.
At zero and one, $\xi$ is nonzero and $H_g$ is entire.
The horizontal bounds hold even near zeros, because
\[
 \Phi(s)=
 \left(\frac1s+\frac1{s-1}-\frac12\log\pi
                       +\frac12\psi(s/2)\right)\zeta(s)M_v(s)
       +\zeta'(s)M_v(s).
\]
The Mellin factor decreases faster than every power, whereas the
other factors have the polynomial or logarithmic bounds proved above.
Cauchy's theorem gives $I_c(\Phi)=I_{1-c}(\Phi)$.
Substitution $s\mapsto1-s$ and \eqref{nwc:eq:logderivative} therefore give
\[
 I_c\bigl(\nu(s)(H_g(s)+H_g(1-s))\bigr)=0.
\]
Equation \eqref{nwc:eq:scalar-contour} proves $L(\mathcal Ev)=0$.
Proposition~\ref{nwc:prop:ideal} and continuity now give
$L(\mathcal J_a)=0$. Taking the infimum over representatives proves
the quotient estimate and compatibility.
\end{proof}

For example, the even Schwartz function
\[
 v_0(x)=x^2\left(x^2-\frac3{2\pi}\right)e^{-\pi x^2}
\]
has $v_0(0)=\int v_0=0$ by the Gaussian second and fourth moments.
Its theta image is a defining relation, so $L(\mathcal Ev_0)=0$.
Every map from $\mathcal Q_{\Theta,a}$ must annihilate that image
as an element, before any scalar evaluation is applied.

\section{Trace-class representations}

Let $\mathfrak S_1(\mathcal H)$ be the ideal of trace-class operators
on a separable positive Hilbert space $\mathcal H$, with trace norm,
ordinary composition, and Hilbert adjoint. A star homomorphism is
complex-linear, multiplicative, and adjoint-preserving. It need not
be injective or nondegenerate. For such a map $\Pi$ put
\[
 \mathcal H_0=\bigcap_{g\in\mathcal A_a}\ker\Pi(g),\qquad
 \mathcal H_{\mathrm e}=\mathcal H_0^\perp.
\]
The second space is its essential representation.

\begin{lemma}\label{nwc:lem:characters}
Every nonzero continuous star character of $\mathcal A_a$ is
$g\mapsto h_g(\lambda)$ for a unique $\lambda\in\mathbb R$.
\end{lemma}
\begin{proof}
Translations satisfy
\begin{equation}\label{nwc:eq:multipliers}
 (R_tf)*g=f*(R_tg)=R_t(f*g),\qquad
 (R_tg)^*=R_{-t}(g^*).
\end{equation}
For a character $\chi$, choose $g_0$ with $\chi(g_0)\ne0$.
The multiplier identities give
\[
 \chi(R_tf)=k(t)\chi(f),\qquad
 k(t)=\frac{\chi(R_tg_0)}{\chi(g_0)}.
\]
Continuity of translations makes $k$ continuous; their group law
gives $k(t+s)=k(t)k(s)$ and $k(0)=1$.
The star identity gives $\overline{k(t)}=k(-t)=k(t)^{-1}$.
A continuous local argument is additive, so $k(t)=e^{it\lambda}$
for a real $\lambda$.

The integral $f*g_0=\int f(t)R_tg_0\,dt$ converges in
$\mathcal A_a$: its $N$th seminorm is bounded by
$p_{a,N}(g_0)\int|f(t)|e^{a|t|}\,dt$.
Applying $\chi$ and dividing by $\chi(g_0)$ proves the formula.
Modulated Gaussians separate frequencies. Real Fourier evaluation
has the required properties, which proves the converse.
\end{proof}

\begin{theorem}\label{nwc:thm:classification}
A continuous star homomorphism
$\Pi:\mathcal A_a\to\mathfrak S_1(\mathcal H)$ has the form
\begin{equation}\label{nwc:eq:diagonal}
 \mathcal H=\mathcal H_0\oplus
     \bigoplus_{\lambda\in\Sigma}\mathbb C^{n_\lambda},\qquad
 \Pi(g)=0\oplus\bigoplus_{\lambda\in\Sigma}h_g(\lambda)I_{n_\lambda},
\end{equation}
where $\Sigma\subset\mathbb R$ is locally finite,
$n_\lambda$ are positive integers, and
$N_\Sigma(R)=\sum_{|\lambda|\leq R}n_\lambda$ has polynomial growth.
The set and its multiplicities are unique.
Conversely, every such counting measure defines a continuous star
homomorphism on its Hilbert direct sum, on every test step.
\end{theorem}
\begin{proof}
Let $q_\epsilon(u)=(4\pi\epsilon)^{-1/2}e^{-u^2/(4\epsilon)}$.
These tests satisfy $q_\epsilon*q_\delta=q_{\epsilon+\delta}$ and
$q_\epsilon*g\to g$ in every $\mathcal A_a$.
For the convergence, use
$q_\epsilon*g-g=\int q_\epsilon(t)(R_tg-g)\,dt$,
translation continuity near zero, and weighted Gaussian tails outside.

Every $\Pi(g)$ is compact and normal, and the images commute with
one another and their adjoints. The compact operator
$K_0=\Pi(q_1)=\Pi(q_{1/2})^*\Pi(q_{1/2})$ is positive and
has kernel $\mathcal H_0$.
Indeed, $K_0v=0$ implies $\Pi(q_{2^{-r}})v=0$ for every $r$ by
repeated square roots, and then
$\Pi(g)v=\lim_r\Pi(g)\Pi(q_{2^{-r}})v=0$.
The compact self-adjoint spectral theorem decomposes
$\mathcal H_{\mathrm e}$ into finite-dimensional nonzero eigenspaces
of $K_0$. Each is invariant under the commuting normal family.
Simultaneous finite-dimensional diagonalization and
Lemma~\ref{nwc:lem:characters} give \eqref{nwc:eq:diagonal}.
The eigenvalue of $K_0$ on a frequency block is $e^{-\lambda^2}$.
Compactness therefore gives local finiteness and finite multiplicities.
The joint eigenspaces give uniqueness.

Continuity supplies $C>0$ and an integer $N$ with
$\|\Pi(g)\|_1\leq Cp_{a,N}(g)$.
For $R\geq1$, set $q_R(u)=Rq_1(Ru)$.
Then $h_{q_R}(\lambda)=e^{-\lambda^2/R^2}$ and
$p_{a,N}(q_R)\leq R^Np_{a,N}(q_1)$.
Consequently
\begin{equation}\label{nwc:eq:counting-necessity}
 N_\Sigma(R)\leq e\|\Pi(q_R)\|_1
              \leq eCp_{a,N}(q_1)R^N.
\end{equation}

Conversely, suppose $N_\Sigma(R)\leq C_0(1+R)^d$.
For an integer $M>d$, dyadic summation gives
$S_M=\sum_\lambda n_\lambda(1+|\lambda|)^{-M}<\infty$.
Integration by parts bounds
$|h_g(\lambda)|\leq2^M(1+|\lambda|)^{-M}p_{0,M}(g)$.
Thus the direct sum is trace class and
\begin{equation}\label{nwc:eq:trace-bound}
 \|\Pi(g)\|_1=\sum_\lambda n_\lambda|h_g(\lambda)|
             \leq2^MS_Mp_{0,M}(g)\leq2^MS_Mp_{a,M}(g).
\end{equation}
Fourier multiplication and real conjugation prove the star identities.
\end{proof}

Its ordinary trace and positive pairing are
\[
 \operatorname{Tr}\Pi(g)=\sum_\lambda n_\lambda h_g(\lambda),\qquad
 \operatorname{Tr}(\Pi(f)^*\Pi(g))
   =\sum_\lambda n_\lambda\overline{h_f(\lambda)}h_g(\lambda).
\]
The multiplicities are dimensions. Rescaling a Hilbert norm does not
change them. Noninteger masses require a different specified trace
weight. The separable direct sum in the theorem can be isometrically
embedded into $\ell^2(\mathbb N_{>0})$, extending the representation
by zero on its orthogonal complement.

\begin{proposition}\label{nwc:prop:generator}
A compatible family into a fixed $\mathfrak S_1(\mathcal H)$ on the
test steps has one common decomposition
\eqref{nwc:eq:diagonal}. On its essential space define
\[
 D=\bigoplus_\lambda\lambda I_{n_\lambda},\qquad
 \operatorname{Dom}D=
 \left\{v:\sum_\lambda\lambda^2\|v_\lambda\|^2<\infty\right\}.
\]
This operator is self-adjoint with compact resolvent.
For $U_t=e^{itD}$,
\begin{equation}\label{nwc:eq:operator-covariance}
 \Pi(R_tg)=U_t\Pi(g),\qquad
 \Pi(g^{(k)})=(-iD)^k\Pi(g).
\end{equation}
The range of every $\Pi(g)$ lies in
$\bigcap_{k\geq0}\operatorname{Dom}D^k$.
If $\|\Pi(g)\|_1\leq Cp_{a,N}(g)$, then
$\|D^k\Pi(g)\|_1\leq Cp_{a,N+k}(g)$.
Any strongly continuous unitary group satisfying the first identity
in \eqref{nwc:eq:operator-covariance} equals this group on the
essential space.
\end{proposition}
\begin{proof}
Compact tests belong to every step and are dense in each.
Their common diagonalization extends by trace-norm limits, with
entries determined by continuous Fourier evaluation.
The diagonal domain gives self-adjointness.
Local finiteness makes the resolvent entries tend to zero, hence the
resolvent is compact.
The transform identities
$h_{R_tg}(\lambda)=e^{it\lambda}h_g(\lambda)$ and
$h_{g^{(k)}}(\lambda)=(-i\lambda)^kh_g(\lambda)$ prove the formulas.
The bound $|\lambda|^k|h_g(\lambda)|\leq\|g^{(k)}\|_1$ proves
the range assertion; trace-norm continuity gives the stated estimate.
On a frequency block, choose $g$ with $h_g(\lambda)\ne0$.
The covariance identity forces the asserted value of $U_t$ there.
\end{proof}

The covariance requirement leaves the unitary action on
$\mathcal H_0$ unspecified; it does not force compact resolvent on
that unused summand. On the essential space the heat operators
$\Pi(q_\epsilon)=e^{-\epsilon D^2}$ tend strongly to the identity.
When that space is infinite-dimensional, its identity is outside
$\mathfrak S_1$.

\begin{corollary}\label{nwc:cor:theta-recipient}
A representation in Theorem~\ref{nwc:thm:classification} descends
to $\mathcal Q_{\Theta,a}$ if and only if
\[
 \Sigma\subseteq\Sigma_0
   :=\{\lambda\in\mathbb R:\zeta(1/2+i\lambda)=0\}.
\]
Its trace-norm estimate descends to the corresponding quotient
seminorm. This condition specifies neither the multiplicities nor
equality of its trace with $L$.
\end{corollary}
\begin{proof}
Descent means that every block vanishes on every theta generator.
Equation \eqref{nwc:eq:real-hull} is exactly that condition.
Continuity annihilates the closed ideal, and taking the infimum over
representatives gives the quotient estimate.
\end{proof}

\section{What equality of the full trace would imply}

We make the multiplicity consequence explicit. Let $Z$ be the set
of distinct zeros of $\xi$, let $m_\rho$ be their multiplicities,
and put $\omega_\rho=(\rho-1/2)/i$ and $m_{\omega_\rho}=m_\rho$.
The following scalar facts involve all zeros, with no reality assumption.

\begin{lemma}\label{nwc:lem:full-scalar}
For every $g\in\mathcal A$,
\begin{equation}\label{nwc:eq:zero-sum}
 L(g)=\sum_{\rho\in Z}m_\rho h_g(\omega_\rho),
\end{equation}
with absolute convergence.
For any finite subset of the distinct ordinates $\omega_\rho$ and
any prescribed values on that subset, there is a test in
$\mathcal A_\infty:=\bigcap_{a>1/2}\mathcal A_a$ taking those
Fourier values and vanishing at every other ordinate.
\end{lemma}
\begin{proof}
The theta integral \eqref{nwc:eq:xi-integral} gives
$\log\max_{|s|\leq R}|\xi(s)|=O(R\log(R+2))$:
bound $\theta(t)$ by $Ce^{-\pi t}$ and the powers by $t^{R+1}$,
then use the gamma integral. Jensen's formula, with $\xi(0)=1/2$,
gives $\sum_{|\rho|\leq R}m_\rho=O(R\log(R+2))$.
The Euler product and functional equation place the zeros in
$0\leq\operatorname{Re}\rho\leq1$.

Here a direct factorization argument suffices. The zero count makes
the canonical product
\[
 P(s)=\prod_{\rho\in Z}\bigl((1-s/\rho)e^{s/\rho}\bigr)^{m_\rho}
\]
locally normally convergent. Choose radii $R_j\in[j,j+1]$ at
distance at least $j^{-2}$ from every zero modulus.
The excluded length is $O(j^{-1}\log j)$, so such radii exist for
large $j$. On $|s|=R_j$, each factor with $|\rho|\leq2R_j$
satisfies $|1-s/\rho|\geq(4R_j^3)^{-1}$.
Partial summation of the zero count gives
\[
 \sum_{|\rho|\leq2R}\frac{m_\rho}{|\rho|}=O(\log^2(R+2)),\qquad
 \sum_{|\rho|>2R}\frac{m_\rho}{|\rho|^2}=O(R^{-1}\log(R+2)).
\]
For the remaining factors use
$|\log((1-z)e^z)|\leq C|z|^2$ when $|z|\leq1/2$.
These estimates give $\log|P(s)|\geq-CR_j\log^2(R_j+2)$
on the selected circles.
The zero-free entire quotient $\xi/P$ is $e^G$ for an entire $G$,
and $\operatorname{Re}G\leq CR_j\log^2(R_j+2)$ there.
If $G(s)=\sum c_ks^k$, integration of its real part on a circle
gives, for $k\geq1$,
\[
 |c_k|\leq 2R_j^{-k}
       \bigl(\max_{|s|=R_j}\operatorname{Re}G(s)
                         -\operatorname{Re}G(0)\bigr).
\]
For $k\geq2$ the right side tends to zero. Thus $G(s)=\alpha+\beta s$,
and logarithmic differentiation gives the locally convergent sum
\[
 \nu(s)=\beta+\sum_{\rho\in Z}m_\rho
             \left(\frac1{s-\rho}+\frac1\rho\right).
\]
Choose heights $T_j\in[j,j+1]$ at distance at least $j^{-2}$
from the imaginary parts of all zeros, and also with the corresponding
property for $-T_j$. Such heights exist for large $j$: the excluded
length is $O(j^{-1}\log j)$ by the zero count.
On the horizontal edges between $1-c$ and $c$, the terms
$|\rho|\leq2T_j$ give $O(T_j^3\log(T_j+2))$.
The remaining terms are bounded by
$O(T_j\sum_{|\rho|>2T_j}m_\rho|\rho|^{-2})=O(\log(T_j+2))$.
The strip decay of $H_g$ therefore makes the horizontal contour
integrals vanish. The residue theorem, followed by
$\nu(1-s)=-\nu(s)$, now identifies
\eqref{nwc:eq:scalar-contour} with the absolutely convergent sum
\eqref{nwc:eq:zero-sum}.

For isolation, set $F(z)=\xi(1/2+iz)/\xi(1/2)$.
Here $\xi(1/2)\ne0$: the alternating series
$\eta(s)=\sum_{n\geq1}(-1)^{n-1}n^{-s}$ converges on
$\operatorname{Re}s>0$ and is positive for real $s>0$, by grouping
successive pairs. The identity
$\zeta(s)=\eta(s)/(1-2^{1-s})$ therefore makes $\zeta(s)<0$
for real $0<s<1$.
The same argument shows that $\xi$ has no real zero.
Consequently every ordinate has nonzero real part, and the ordinates
occur in pairs $\omega,-\omega$.
Applying the same factorization argument to $F$ and pairing its
opposite zeros gives
\begin{equation}\label{nwc:eq:paired-product}
 F(z)=\prod_{\operatorname{Re}\omega>0}
                         (1-z^2/\omega^2)^{m_\omega}.
\end{equation}
It converges locally normally since
$\sum m_\omega|\omega|^{-2}<\infty$.
The quotient of $F$ by this product is an exponential of a linear
polynomial; evenness and $F(0)=1$ make it one.
Integration of the zero-counting bound gives a bound
$\exp(C(1+|z|)\log(2+|z|))$ for the product.

For a finite set $E$ of ordinates, remove their full factors:
$F_E(z)=F(z)/\prod_{\omega\in E}(z-\omega)^{m_\omega}$.
It is nonzero at each point of $E$ and has the same type of growth
bound. For any $t>0$, Lagrange interpolation supplies a polynomial
$P$ for which $G(z)=e^{-tz^2}F_E(z)P(z)$ has the prescribed values
on $E$. It vanishes at every other ordinate.
On every horizontal strip, $G$ and all its derivatives decrease
rapidly. Define $g(u)=(2\pi)^{-1}\int G(r)e^{-iru}\,dr$.
For $u>0$ move the contour to $\operatorname{Im}z=-b$, and for
$u<0$ to $\operatorname{Im}z=b$.
Applying the same shifts to $z^kG(z)$ gives
$|g^{(k)}(u)|\leq C_{b,k}e^{-b|u|}$ for every $b>0$.
Thus $g\in\mathcal A_\infty$ and $h_g=G$, as required.
\end{proof}

For a compact self-adjoint Hilbert--Schmidt operator $K$, define
\[
 \det\nolimits_2(I-zK)
   =\prod_{\kappa\ne0}(1-z\kappa)e^{z\kappa},
\]
where the nonzero eigenvalues $\kappa$ are repeated with their
multiplicities. Square summability makes this product locally normally
convergent.

\begin{theorem}[Conditional full trace consequences]
\label{nwc:thm:conditional}
Suppose a compatible theta-quotient star representation satisfies
\begin{equation}\label{nwc:eq:full-comparison}
 \operatorname{Tr}\Pi(g)=L(g)\qquad(g\in\mathcal A).
\end{equation}
Then every ordinate $\omega_\rho$ is real, every one occurs in
\eqref{nwc:eq:diagonal}, and its Hilbert multiplicity is $m_\rho$.
The essential generator is invertible. Its inverse $K=D^{-1}$ is
bounded, self-adjoint, and Hilbert--Schmidt.
For $F$ as in \eqref{nwc:eq:paired-product}, define
\[
 t_m=-\frac{(\log F)^{(m)}(0)}{(m-1)!},\qquad m\geq2.
\]
Then
\begin{equation}\label{nwc:eq:all-moments}
 \operatorname{Tr}K^m=t_m\quad(m\geq2),\qquad
 \det\nolimits_2(I-zK)=F(z).
\end{equation}
\end{theorem}
\begin{proof}
By Corollary~\ref{nwc:cor:theta-recipient}, all represented
frequencies lie in $\Sigma_0$.
For a real ordinate $\gamma$, apply
Lemma~\ref{nwc:lem:full-scalar} to the test whose Fourier values
are one at $\gamma$ and zero at every other ordinate.
Its operator is the projection onto that block, or zero if the
block is absent. Equality of traces forces $n_\gamma=m_\gamma$.
If a nonreal ordinate existed, its isolating test would have zero
operator but scalar value $m_\omega>0$, a contradiction.

There is no ordinate zero, and local finiteness separates the
essential spectrum from zero. The zero count gives
$\sum_\gamma m_\gamma|\gamma|^{-2}<\infty$.
Hence $K$ has the stated properties and every $K^m$, $m\geq2$,
is trace class. Taking local logarithms in the paired product gives
$t_m=\sum_\gamma m_\gamma\gamma^{-m}=\operatorname{Tr}K^m$.
For the regularized determinant, pairing the factors
$(1-z/\gamma)e^{z/\gamma}$ cancels their exponential terms and
gives precisely \eqref{nwc:eq:paired-product}.
\end{proof}

The cyclic moment measure is
$\sum_\gamma m_\gamma\gamma^{-2}\delta_{\gamma^{-1}}$;
the counting representation has $m_\gamma$ dimensions at $\gamma$.
Their comparison can be made in the Hilbert--Schmidt space
$\mathfrak S_2(\mathcal H_{\mathrm e})$.
With $a_j=K^{j+1}$, all $a_j$ are Hilbert--Schmidt and
\[
 a_i^*a_j\in\mathfrak S_1,
 \qquad \langle a_i,a_j\rangle_{\mathfrak S_2}=t_{i+j+2}.
\]
Left multiplication by $K$ is a bounded self-adjoint operator there;
its cyclic subspace generated by $K$ has the stated moment measure.
In general $a_0=K$ need not be trace class. These are conditional
states in the square-integrable domain of the trace, with integrable
products. The full comparison \eqref{nwc:eq:full-comparison} has not
been constructed from lattice operations.

\section{Coefficient complexes and supported intersections}

The relation ideal gives a two-term coefficient dg algebra
\begin{equation}\label{nwc:eq:coefficient-dg}
 \mathcal C_{\Theta,a}^{-1}=\mathcal J_a,
 \qquad \mathcal C_{\Theta,a}^0=\mathcal A_a,
 \qquad d:\mathcal J_a\hookrightarrow\mathcal A_a.
\end{equation}
Products in degree zero are convolution, mixed products are the ideal
action, and the product of two degree-minus-one elements is zero.
Use the graded involution convention
\[
 (xy)^*=(-1)^{|x||y|}y^*x^*,\qquad d(x^*)=(dx)^*.
\]
The inherited involution satisfies these identities.
The differential has degree $+1$.
The only nontrivial odd Leibniz identity is
$0=d(jk)=(dj)k-j(dk)$, whose two terms are the same convolution.
Thus this is a continuous graded-commutative dg algebra, with
$H^{-1}=0$ and $H^0=\mathcal Q_{\Theta,a}$.
Extend $L$ by zero in degree minus one. Since $L(dj)=0$, this
extension is a closed Hermitian degree-zero trace; its positivity
is not asserted.

On its total coefficient space use
$q_{a,N}(g,j)=p_{a,N}(g)+p_{a,N}(j)$.
Multiplication and the differential have seminorm bounds one.
A quotient map $\Pi$ to an algebra in degree zero is a chain map
from \eqref{nwc:eq:coefficient-dg} exactly when it kills
$\mathcal J_a$.
This permits a precise coefficient extension on finite supported
intersections, conditional on that map.

Let $G$ be a finite directed graph with real edge gains $a_e$,
and let $X_G=\mathbb C^{V(G)}$ with coordinates $u_v$.
Write $s(e),t(e)$ for the source and target vertices of an edge, and put
\[
 \ell_e=u_{t(e)}-u_{s(e)}-a_e,
 \qquad
 \mathcal K_G=\mathcal O_{X_G}\otimes
              \Lambda(\eta_e:e\in E(G)),
 \quad |\eta_e|=-1,
 \quad d\eta_e=\ell_e.
\]
This is the Koszul complex of the full list of edge equations.
It represents their derived zero locus, retaining dependencies
between the equations. Away from their common zero set, a nonvanishing
$\ell_e$ gives the contraction $\ell_e^{-1}\eta_e\wedge-$.

For a complete Fr\'echet dg coefficient algebra $E$, let
$\mathcal O(U;E)$ be the space of holomorphic $E$-valued functions
with its compact-open coefficient seminorms. Locally on polydiscs
this is $\mathcal O(U)\widehat\otimes_\pi E$.
One obtains the identification by the holomorphic power series:
Cauchy bounds make the coefficient series converge in each projective
seminorm on a smaller polydisc. These local spaces form a sheaf.
Put
\[
 \mathcal K_G(E)(U)=\mathcal O(U;E)\otimes\Lambda(\eta_e),
 \qquad
 d(x\otimes\eta)=d_Ex\otimes\eta+(-1)^{|x|}x\otimes d\eta.
\]
The exterior factor is finite-dimensional. Its displayed tensor
products therefore require no further completion.

\begin{proposition}\label{nwc:prop:supported-coefficients}
Every continuous dg algebra map $T:E\to E'$ induces a continuous
map $\mathcal K_G(E)\to\mathcal K_G(E')$.
It commutes with affine pullbacks that preserve the full list of
edge equations, including equations that become zero or nonzero
constants after a vertex collision.
For a balanced triangle and a closed degree-zero coefficient $x$, it carries the cycle
$x(\eta_{13}-\eta_{12}-\eta_{23})$ to
$T(x)(\eta_{13}-\eta_{12}-\eta_{23})$.
For $E=\mathcal C_{\Theta,a}$ this yields an injection
\begin{equation}\label{nwc:eq:triangle-injection}
 \mathcal Q_{\Theta,a}\longrightarrow
 H^{-1}\bigl(\mathcal K_\Delta(E)(\mathbb C^3)\bigr).
\end{equation}
If $E'$ is concentrated in degree zero, the displayed image of a
cycle is nonzero in cohomology whenever $T(x)\ne0$.
\end{proposition}
\begin{proof}
Apply $T$ to every coefficient. The projective tensor property and
the finite exterior factor prove continuity; the formula for $d$
proves the chain identity.
To describe collisions explicitly, let $p:V(G)\twoheadrightarrow W$
and put $u_v=y_{p(v)}+b_v$.
Then
\[
 \ell_e\longmapsto
 y_{p(t(e))}-y_{p(s(e))}
       -(a_e+b_{s(e)}-b_{t(e)}).
\]
Keep the edge label and its exterior generator even when its endpoints
coincide. This substitution defines the pullback Koszul map and
commutes with $T$. Two substitutions compose by substitution of their
affine coordinate formulas; their dg maps consequently agree exactly.
No ordering or consecutiveness of the identified vertices is used.
If a loop equation is a nonzero constant, the resulting complex is
contractible. A zero equation retains its exterior generator.

For the triangle, take $a_{13}=a_{12}+a_{23}$ and introduce
$w=u_1$, $x_1=\ell_{12}$, $x_2=\ell_{23}$ and
$\kappa=\eta_{13}-\eta_{12}-\eta_{23}$.
Then $d\eta_{12}=x_1$, $d\eta_{23}=x_2$, and $d\kappa=0$.
Evaluation at $w=x_1=x_2=0$, with
$\eta_{12},\eta_{23}$ sent to zero, is a chain retraction to
$E\otimes\Lambda(\kappa)$.
For $E=\mathcal C_{\Theta,a}$, every degree-zero coefficient $x$
gives a cycle $x\kappa$.
Changing $x$ by $j\in\mathcal J_a$ changes that cycle by
$d(j\kappa)$, with $j$ here in degree minus one.
The retraction shows that it can be a boundary only when
$x\in\mathcal J_a$. This proves \eqref{nwc:eq:triangle-injection}.
The same retraction in a degree-zero target proves the last assertion.
\end{proof}

In particular, gains $\log2,\log3,\log6$ retain a degree-minus-one
class. A coefficient map to degree-zero operators extends to the
whole Koszul complex by leaving its exterior factor present.
Collapsing the complex to the scalar trace would lose this class.

To specify a finite derived-Hom calculation, adjoin a scalar unit,
$E^+=\mathbb C\oplus E$, and put
\[
 B_E(U)=\mathcal O(U;E^+),\qquad
 P_G(E)(U)=B_E(U)\otimes\Lambda(\eta_e).
\]
Use the same total differential as above. The unitized module
$P_G(E)(U)$ is finite free as a graded $B_E(U)$-module.
Its proper coefficient ideal
$I_G(E)(U)=\mathcal K_G(E)(U)$ need not be free.
Between two unitized modules $P,Q$ use the finite matrix complex
\begin{equation}\label{nwc:eq:continuous-hom}
 \operatorname{Hom}_{B_E(U),\mathrm{cont}}^\bullet(P,Q),
 \qquad d\phi=d_Q\phi-(-1)^{|\phi|}\phi d_P.
\end{equation}
Its topology is the corresponding finite product of coefficient
Fr\'echet spaces.
The increasing filtration by exterior degree is finite, is preserved
by the differential, and has successive quotients that are finite sums
of shifts of $B_E(U)$. Since continuous Hom from $B_E(U)$ is the
identity functor, Hom from $P$ preserves quasi-isomorphisms by this
filtration. Hence the displayed complex computes derived Hom for
these unitized modules.

The distinction matters even on a point with no edges.
Take $E=\mathbb C\varepsilon$ with $\varepsilon^2=0$ and zero
differential. Then $B_E=\mathbb C[\varepsilon]/(\varepsilon^2)$,
whereas $I_G(E)=E$. The exact resolution
\[
 \cdots\xrightarrow{\varepsilon}B_E
       \xrightarrow{\varepsilon}B_E
       \xrightarrow{\varepsilon}E\longrightarrow0
\]
shows that $\operatorname{Ext}_{B_E}^n(E,E)=\mathbb C$ for every
$n\geq0$: applying $\operatorname{Hom}_{B_E}(-,E)$ makes each
differential zero. Its degree-zero Hom complex alone would miss all
the positive Ext groups. Thus the finite matrix calculation does not
compute derived Hom for the proper coefficient ideals, or for a
chiral field recipient. No finite trace is assigned to the adjoined unit.

Normal differentiation already distinguishes the two settings.
For a constant coordinate vector $v$, let $\partial_v$ differentiate
the holomorphic coefficient and fix every $\eta_e$. Directly,
\begin{equation}\label{nwc:eq:normal-defect}
 [\partial_v,d]
       =\sum_e(\partial_v\ell_e)\,\iota_e,
\end{equation}
where $d$ is the total differential and $\iota_e$ is total contraction:
\[
 \iota_e(x\otimes\eta)=(-1)^{|x|}x\otimes\iota_e^{\Lambda}(\eta),
 \qquad \iota_e^{\Lambda}(\eta_f)=\delta_{ef}.
\]
Thus a normal derivative is generally not a chain endomorphism of
this complex. It needs an appropriate jet or differential-module
recipient if it is to be part of a supported comparison.
The finite Koszul construction alone does not provide that recipient
or the required field operations.

\begin{proposition}\label{nwc:prop:derivation-distinction}
Suppose $\Pi:\mathcal A_a\to\mathfrak S_1(\mathcal H)$ is a
continuous star homomorphism and a derivation $T$ of its image
satisfies $T\Pi(g)=\Pi(g')$ for every $g$.
Then its only possible frequency is zero.
If it also descends through $\mathcal Q_{\Theta,a}$, it is zero.
\end{proposition}
\begin{proof}
Differentiation under convolution gives
$(f*g)'=f'*g=f*g'$.
The derivation rule for $T$, compared with multiplicativity of $\Pi$,
therefore gives
\[
 \Pi((f*g)')=\Pi(f'*g)+\Pi(f*g')=2\Pi((f*g)').
\]
The Gaussian approximate identity implies $\Pi(g')=0$ for every $g$.
The decomposition \eqref{nwc:eq:diagonal} then forces
$\lambda h_g(\lambda)=0$ for all tests, hence $\lambda=0$ on every
block. But $\zeta(1/2)\ne0$, so Corollary~\ref{nwc:cor:theta-recipient}
excludes that block for a theta-quotient representation.
\end{proof}

This proposition concerns the specific identification of a convolution
multiplier with a derivation of the same product. A supported lift
can instead keep native dilation as a module operation, as in
\eqref{nwc:eq:operator-covariance}. Its geometric translation,
normal derivatives, and collision operations still require their
own comparison maps.

\section{Tensor control and the remaining positive comparison}

For completed projective coefficient tensors, a bound
$\|\Pi(g)\|_1\leq Cp(g)$ implies
\[
 \|\Pi^{\widehat\otimes r}(x_r)\|_\pi
       \leq C^r p_\pi^{(r)}(x_r).
\]
For $r\geq1$, multiplication of $r$ trace-class images has norm
at most one, since
$\|AB\|_1\leq\|A\|_{\mathrm{op}}\|B\|_1\leq\|A\|_1\|B\|_1$.
On the nonfactorial entire tensor space the bound is consequently
\begin{equation}\label{nwc:eq:entire-bound}
 \sum_{r\geq0}R^r\|\Pi^{\widehat\otimes r}(x_r)\|_\pi
 \leq\sum_{r\geq0}(CR)^r p_\pi^{(r)}(x_r),\qquad R>0.
\end{equation}
The term $r=0$ is interpreted in the scalar coaugmentation.
It is not assigned a finite ordinary operator trace.

For \eqref{nwc:eq:coefficient-dg}, the internal tensor differential
has at most $r$ terms at length $r$, and the adjacent multiplication
part of the reduced bar differential has at most $r-1$ terms.
From radius $\lambda R$ to $R$, $\lambda>1$, their bounds are
\[
 \sup_{r\geq0}r\lambda^{-r},\qquad
 R^{-1}\sup_{r\geq2}(r-1)\lambda^{-r}.
\]
Simultaneous translation of $r$ factors has coefficient bound
$e^{ar|t|}$ and is continuous from radius $Re^{a|t|}$ to $R$.
This is a tensor-space bound; simultaneous translation is not claimed
to commute with the bar differential. Indeed,
$(R_tf)*(R_tg)=R_{2t}(f*g)$.
The factorial norms with weights $R^r/r!$ define a different
completion. A change of bar differential signs by $(-1)^r$ does
not identify these topologies.

The exponential weight $a$ controls arithmetic translation in the
coefficient variable. The radius $R$ controls tensor length.
Neither estimate supplies the other, and the coefficient tensor map
does not by itself construct a field collision substitution.
The scalar empty tensor, a local field vacuum, and a cyclic Hilbert
vector are distinct data. No local field vacuum has been chosen here.

A positive compression illustrates a remaining algebraic condition.
Suppose $U_t$ is a strongly continuous unitary group on a Hilbert
space, $T(g)=\int g(t)U_t\,dt$ is its strong operator integral,
and $P$ is an orthogonal projection. Set $T_P(g)=PT(g)P$ on
$P\mathcal H$. Its exact multiplicativity defect is
\[
 T_P(f*g)-T_P(f)T_P(g)=PT(f)(1-P)T(g)P.
\]
For $g=f$ and the first input $f^*$, this defect is
$((1-P)T(f)P)^*((1-P)T(f)P)$.
Multiplicativity on all adjoint squares therefore forces $P$ to
reduce every $T(f)$, and then every $U_t$ by shifted approximate
identities. Positivity of the compression alone does not supply
this reduction property or the full trace identity.

An independent positive comparison would follow from a geometric
Hilbert space with unitary dilation, a compact-resolvent generator
with polynomial eigenvalue count, operator annihilation of every
theta relation, and equality of its trace with \eqref{nwc:eq:weil}
on the compact core. The preceding estimates would extend that
equality to every test step and give its coefficient derivative and
tensor domains. Constructing this Hilbert space and proving its full
trace identity remain open. Prescribing zero ordinates or discarding
possible nonreal characters supplies neither proof.

For a scalar illustration of the prime-square term, take
$\delta=1/64$ and let
\[
 \chi(u)=c\exp\!\left(-\frac1{1-(u/\delta)^2}\right)
 \quad(|u|<\delta),\qquad \chi(u)=0\quad(|u|\geq\delta),
 \qquad \int\chi(u)^2\,du=1.
\]
Put $\ell=\log4$, $f=\chi+R_\ell\chi$, and $g=f^**f$.
Its support consists of intervals of radius $2\delta$ about
$0,\ell,-\ell$, and $g(\ell)=g(-\ell)=1$.
Among the points $\pm\log n$, $n\geq2$, only $\pm\log4$ occur
in these intervals. Thus the entire prime contribution to $L(g)$
is $-\log2$. Removing those two contributions from that scalar
expression changes it by $+\log2$; the polar and archimedean terms
remain present. An operator comparison on this input would also
have to prove $\Pi(g)=\Pi(f)^*\Pi(f)$ and identify the two
geometric trace contributions before omitting them.

The full supported comparison additionally requires its field or sheaf
recipient, pull-push operations, all collision homotopies, normal
derivatives, Fourier and dilation compatibility, and the local vacuum
on a common integrable domain. The finite Koszul extension identifies
one supported coefficient structure to preserve. It proves neither
that full comparison nor a universal obstruction to it.

